\subsubsection*{Thesis title}  
\Title

\subsubsection*{Abstract} 
The objective of this thesis was to develop and evaluate a reinforcement-learning approach for training a computer-controlled player in the logic-based board game Rummikub. Due to the large state and action spaces and the presence of sparse and delayed rewards, a Curriculum Learning environment was designed to introduce selected game mechanics through tasks of gradually increasing difficulty. The environment and training framework were implemented using Unity and the ML-Agents Toolkit.

The experimental study compared two configurations: Proximal Policy Optimization (PPO) as a baseline and PPO supported by Generative Adversarial Imitation Learning. Ten independent agents were trained for each configuration. Their performance was assessed using curriculum progression, the number of environment steps, wall-clock training time, and the win rate obtained over 1,000 evaluation episodes per agent.

\subsubsection*{Key words}  
Reinforcement learning, imitation learning, Curriculum Learning, PPO, GAIL, Unity ML-Agents, Rummikub

\subsubsection*{Tytuł pracy}
\begin{otherlanguage}{polish}
\TitleAlt
\end{otherlanguage}

\subsubsection*{Streszczenie} 
\begin{otherlanguage}{polish}
Celem niniejszej pracy było opracowanie i ocena podejścia opartego na uczeniu się przez wzmocnienie, służącego do szkolenia komputerowego gracza w opartej na logice grze planszowej Rummikub. Ze względu na rozległe przestrzenie stanów i działań oraz występowanie rzadkich i opóźnionych nagród zaprojektowano środowisko uczenia się oparte na programie nauczania (Curriculum Learning), mające na celu wprowadzanie wybranych mechanizmów gry poprzez zadania o stopniowo rosnącym stopniu trudności. Środowisko oraz framework szkoleniowy zaimplementowano przy użyciu silnika Unity oraz zestawu narzędzi ML-Agents Toolkit.

W badaniu eksperymentalnym porównano dwie konfiguracje, Proximal Policy Optimization (PPO) jako punkt odniesienia oraz PPO wspomagane metodą Generative Adversarial Imitation Learning (GAIL). Dla każdej konfiguracji wyszkolono dziesięciu niezależnych agentów. Ich wyniki oceniono na podstawie postępów w programie nauczania, liczby kroków w środowisku, czasu szkolenia w czasie rzeczywistym oraz wskaźnika zwycięstw uzyskanego w 1 000 epizodów ewaluacyjnych na agenta.
\end{otherlanguage}

\subsubsection*{Słowa kluczowe} 
\begin{otherlanguage}{polish}
Uczenie ze wzmocnieniem, uczenie przez naśladowanie, Curriculum Learning, PPO, GAIL, Unity ML-Agents, Rummikub
\end{otherlanguage}

